\subsection*{Materiales}
\begin{itemize}
    \item Matraz Erlenmeyer
        \item Fuente de calentamiento (estufa o placa)
	    \item Termómetro
	        \item Soporte universal
		\end{itemize}
		\subsection*{Procedimiento experimental}
		\begin{enumerate}
		    \item Se colocaron 350 mL de agua en el matraz Erlenmeyer y se elevó su temperatura hasta el rango de 75\,°C a 100\,°C. Posteriormente, se desconectó o retiró la fuente de calor.
		        \item Se midió la temperatura inicial \(T_0\) del agua caliente y se anotó la temperatura ambiente del laboratorio.
			    \item Se registró el descenso de la temperatura en función del tiempo. Esto se realizó con un cronómetro para la toma manual de datos. Se conservaron los valores registrados manualmente.
			        \item Debido a que el enfriamiento fue más rápido al comienzo, se efectuaron las primeras mediciones con intervalos de tiempo cortos. A medida que la temperatura disminuyó más lentamente, se alargaron gradualmente dichos intervalos.
				    \item Se finalizó la adquisición de datos una vez que hubo transcurrido el tiempo total \(T\) establecido para el experimento.
				    \end{enumerate}
